\section{Deployment and Operation}

\subsection{Installation}

The entire Python environment is project-local. Nothing is installed
system-wide.

\begin{lstlisting}
python3.11 -m venv .venv
./.venv/bin/pip install torch --index-url https://download.pytorch.org/whl/cu128
./.venv/bin/pip install -r requirements.txt
./scripts/01_pull_models.sh          # approximately 24 GB of weights
\end{lstlisting}

\noindent The CUDA~12.8 index is mandatory on Blackwell: earlier PyTorch
binaries lack \code{sm\_120} kernels and fail at runtime with
\code{no kernel image is available for execution}.

\subsection{Operating commands}

\begin{center}
\renewcommand{\arraystretch}{1.2}
\begin{tabularx}{\linewidth}{@{}l X@{}}
\toprule
\textbf{Command} & \textbf{Purpose} \\
\midrule
\code{make check} & Preflight: GPU, models, sandbox capability, binaries \\
\code{make index} & Ingest \code{data/corpus} into vectors, graph and communities \\
\code{make watch} & Live indexing dashboard with rate and estimated completion \\
\code{make serve} & API and interface on \code{127.0.0.1:8077} \\
\code{make demo} & Scripted run covering all five rubric points \\
\code{make test} & 61 tests \\
\code{make lockdown} & Install nftables egress denial (requires root) \\
\bottomrule
\end{tabularx}
\end{center}

\subsection{Document placement}

\begin{center}
\renewcommand{\arraystretch}{1.2}
\begin{tabularx}{\linewidth}{@{}l X@{}}
\toprule
\textbf{Directory} & \textbf{Contents} \\
\midrule
\code{data/corpus/} & \textbf{Source documents.} \code{.pdf}, \code{.txt},
  \code{.md}; subdirectories are traversed \\
\code{data/workspace/} & Agent scratch space; file tools are jailed here \\
\code{data/artifacts/} & Generated deliverables \\
\code{data/stores/} & Graph, vectors, extraction cache, audit log \\
\bottomrule
\end{tabularx}
\end{center}

\noindent The filename becomes the \code{doc\_id} used in every citation, so
documents should be named with their real identifiers.
Images (\code{.png}, \code{.jpg}, \code{.tiff}) are \emph{not} automatically
indexed --- they are readable on demand through the vision tools, or may be
converted to PDF for indexing.

\subsection{Air-gap procedure}

\begin{lstlisting}
tar czf models.tar.gz -C ~/.ollama models        # model weights
./.venv/bin/pip download -r requirements.txt -d ./wheels
export HF_HUB_OFFLINE=1 TRANSFORMERS_OFFLINE=1
sudo ./scripts/04_lockdown.sh on                 # egress denial
\end{lstlisting}

\noindent The offline rehearsal should be performed in the first week of a
deployment, not the last. A dependency that resolves a hostname on import is
common and is only discovered by testing with the interface down.

\subsection{Operational cautions}

\begin{description}[leftmargin=0pt,style=nextline]
\item[Single writer.] Indexing and serving cannot run concurrently; the graph
store holds an exclusive lock. Stop the server before re-indexing.
\item[Schema changes require a rebuild.] Altering the enabled ontology packs
changes the database schema. Delete \code{data/stores/graph} and re-index.
\item[The canary reports honestly.] On a machine with ordinary internet access
it will correctly report that an external call succeeded. That is the monitor
working, not failing. Apply the lockdown script to see it blocked.
\end{description}
